% Reciprocal finite proofs; not a transcription of proofs omitted by Carnap's appendix.
\section{Carnap: finite constructions and bounded historical results}
\label{sec:annex-p15-carnap-logical-probability-and-induction}

The language, ranges, and widths are those of
Chapter~\ref{sec:p15-carnap-logical-probability-and-induction}. Unless
specified otherwise all domains and sums are finite. These proofs
reconstruct the results stated in the reader; source locators identify
the corresponding original claims, not an assertion that these exact
proofs occur there.

\begin{annexitem}{annex:p15.carnap-language-counts}{thm:p15-carnap-language-counts}{Profiles, structures, and their cardinalities}
\begin{proof}
Each assignment of \(q\) truth values determines exactly one
\(Q\)-conjunction, and different assignments disagree on a literal.
Thus the \(2^q=k\) conjunctions are mutually exclusive and exhaustive.
Evaluate a molecular property on each of these profiles; its true
profiles determine its union uniquely. Specifying \(d\) independent
literals leaves \(q-d\) truth values free, giving \(2^{q-d}\)
profiles and the claimed width, including \(d=0\).

A state chooses one profile independently for each named individual,
so there are \(k^N\) syntactically possible states. This independence
is freedom of truth-value assignment, not stochastic independence under
an arbitrary measure. Individual permutations preserve counts. Conversely,
if two words have the same counts, match their positions category by
category; the resulting bijection of positions permutes one word into
the other. Consequently the count vectors are precisely the orbits.

The nonnegative vectors summing to \(N\) are in bijection with strings
of \(N\) stars and \(k-1\) separators: counts record the stars
between successive separators, allowing empty blocks. There are
\(\binom{N+k-1}{k-1}\) such strings. For fixed counts, choose the
positions of category one, then category two among the remaining
positions, and so forth. The product of these binomial choices telescopes
to \(N!/\prod_i n_i!\). Zero counts use \(0!=1\).
\end{proof}
\end{annexitem}

\begin{annexitem}{annex:p15.carnap-regular-laws}{thm:p15-carnap-regular-laws}{Regular confirmation and symmetry}
\begin{proof}
Since each state has positive mass, a range has positive mass exactly when
it is nonempty. Summing masses gives a normalized additive measure on
the subsets of the state space. Truth conditions give
\(R(h\land j)=R(h)\cap R(j)\),
\(R(\neg h)=\Omega_N\setminus R(h)\), and
\(R(h\lor j)=R(h)\cup R(j)\). The subset
\(R(e\land h)\subseteq R(e)\) establishes the bounds on the
quotient; equal ranges give equal numerators or denominators. If
\(R(e)\subseteq R(h)\), the numerator equals the denominator;
if \(R(e)\cap R(h)=\varnothing\), it equals zero.

Partitioning \(R(e)\) by \(h\) and \(\neg h\) gives the
complement identity. Inclusion--exclusion for the two numerator ranges
gives the union identity. If \(R(e\land h)\ne\varnothing\),
\[
 \frac{m(e\land h\land j)}{m(e)}
 =\frac{m(e\land h)}{m(e)}
   \frac{m(e\land h\land j)}{m(e\land h)},
\]
which is the multiplication identity. The hypothesis on its denominator
cannot be suppressed.

For a name permutation \(\pi\), satisfaction is preserved by
simultaneously permuting the state and sentence. This follows first for
ground atoms and identity, then for connectives; for a quantifier its
instances are bijectively reindexed by \(\pi\). Hence
\(R(\pi h)=\pi R(h)\). Corresponding states are isomorphic and
have the same mass under a symmetric measure. Summing over this bijection
proves measure invariance, and applying it to numerator and denominator
proves confirmation invariance. The next annex item exhibits two
symmetric measures with different predictions, proving nonuniqueness.
These are finite reconstructions of D55-1--4, T59-1, and T91-2
\cite{Carnap1950Probability}.
\end{proof}
\end{annexitem}

\begin{annexitem}{annex:p15.carnap-star-construction}{thm:p15-carnap-star-construction}{Equal structure weights, projectivity, and equal-state weights}
\begin{proof}
The finite orbits partition the state space. Each orbit \(O\) receives
\(|O|/(r_N|O|)=1/r_N\), so total mass is one and every state has
positive mass. The mass is constant on an orbit and therefore symmetric.
Conversely, a symmetric measure is constant on each orbit. If all orbit
masses are equal, normalization forces each to be \(1/r_N\), and
the mass of each state in \(O\) must be \(1/(r_N|O|)\).
This argument also applies to a finite language with relations.

In the monadic case substitute the cardinalities of the preceding
counting theorem:
\[
 \frac1{\binom{N+k-1}{k-1}\,N!/\prod_i n_i!}
 =\frac{(k-1)!\prod_i n_i!}{(N+k-1)!}.
\]
Let \(z\) be a length-\(N\) word with counts \(n_i\).
Appending category \(i\) gives a word \(zi\) for which
\[
 \frac{m_{N+1}^*(zi)}{m_N^*(z)}=\frac{n_i+1}{N+k}.
\]
Summing over the \(k\) possible categories gives one. Thus the
extension masses sum to \(m_N^*(z)\). Iteration proves the fitting
property for any larger domain and any event determined by the first
\(N\) profiles. Individual symmetry allows any fixed finite list in
place of this initial list. Summing over its range proves the assertion
about sentences using just that list without domain-wide quantification.

The uniform-state measure is \(k^{-N}\) at each word. For any
specified word on \(s\) other individuals, each profile for a fresh
individual has \(k^{N-s-1}\) extensions and hence the same conditional
mass \(1/k\). For a union of such evidence words, the numerator for
a fresh profile is still \(1/k\) of the denominator by finite
additivity. A property consisting of \(w\) profiles has conditional
mass \(w/k\). Both constructions are symmetric, but in the three-name,
binary example \(m^*\) assigns masses \(1/4\) to each constant
word and \(1/12\) to every other word. The two extensions of
\(P(a_1)P(a_2)\) therefore have masses \(1/4,1/12\), giving
\(3/4\) rather than \(1/2\). This also proves the contrast invoked
in the preceding item. The original construction is sec.~110A,
equations (1)--(4), pp.~563--566 \cite{Carnap1950Probability}.
\end{proof}
\end{annexitem}

\begin{annexitem}{annex:p15.carnap-predictive}{thm:p15-carnap-predictive}{Grouping profiles, prediction, and frequency estimates}
\begin{proof}
By fitting, it suffices to consider a specified list of \(s\)
individuals as the whole finite domain for computing its cylinder mass.
Let its coarse category counts be \(s_j\), and let the \(j\)-th
category consist of \(w_j\) complete profiles. Sum the full-profile
mass over refinements of this specified coarse word. Within category
\(j\), a refinement count vector \((v_1,\ldots,v_{w_j})\)
summing to \(s_j\) has \(s_j!/\prod_i v_i!\) assignments to
its fixed positions. Multiplying this number by \(\prod_i v_i!\)
from the full-profile mass leaves \(s_j!\), independent of the
refinement counts. There are \(\binom{s_j+w_j-1}{w_j-1}\)
such counts. Hence the mass of the specified coarse word is
\[
 g(s_1,\ldots,s_p)
 =\frac{(k-1)!}{(s+k-1)!}
   \prod_j\frac{(s_j+w_j-1)!}{(w_j-1)!}.
\]
For the empty list take \(g=1\); the displayed factorial expression
has that value as well. This derivation makes no assumption that coarse
categories themselves are equiprobable structures.

For a specified further coarse word with counts \(t_j\), divide
\(g(s_1+t_1,\ldots,s_p+t_p)\) by \(g(s_1,\ldots,s_p)\).
There are \(t!/\prod_j t_j!\) such further words in the count
report, all with the same conditional mass. Their disjoint union has
conditional mass
\[
 \frac{t!}{\prod_j t_j!}
 \frac{(s+k-1)!}{(s+t+k-1)!}
 \prod_j\frac{(s_j+t_j+w_j-1)!}{(s_j+w_j-1)!},
\]
which is precisely the binomial-coefficient formula in the reader.
If evidence says only its counts, it is a disjoint union of
\(s!/\prod_j s_j!\) specified evidence words. The future event
has the same conditional probability for each, because the future names
are disjoint and permutations of the evidence names leave it fixed.
The common factor therefore cancels from numerator and denominator.

Taking \(t=1\) gives \((s_j+w_j)/(s+k)\). If full profiles are
known, sum \((s_i+1)/(s+k)\) over the profiles constituting \(M\);
the sum is \((s_M+w)/(s+k)\). Thus extra profile information does
not change this particular singular prediction. For \(s>0\), split
the numerator into \(s_M+w\) to obtain the convex-combination
identity. Its nonnegative weights sum to one; as \(s\) increases
with fixed vocabulary, the prior weight \(k/(s+k)\) tends to zero.
If \(s_M/s\to r\), the prediction tends to \(r\).

Write the unobserved count as a sum of indicators of membership in
\(M\). Every indicator has the just-computed conditional mean.
Finite additivity of expectation gives expected unobserved count
\((N-s)(s_M+w)/(s+k)\), and division by the number of unobserved
individuals gives its relative-frequency estimate when that number is
positive. Adding the already observed count gives the whole-population
estimate, valid also for \(N=s\). The inverse distribution follows
by setting the future list to the remainder of the domain; impossible
counts correspond to empty events and have mass zero.

In the reader's two-predicate example, \(P\) contains two profiles,
\(PR\) one, and \(P\lor R\) three. The corresponding observed
counts are \(8,6,9\); substitution gives \(10/14,7/14,12/14\).
For one positive and one negative future \(P\)-instance, substitution
gives \(10\cdot4/105=8/21\). For three observed positives, the
one-predicate value is \(4/5\); the two-predicate value is \(5/7\).
These comparisons change the vocabulary rather than the individual-domain
extension governed by projectivity. The source formulas are sec.~110C,
equations (6)--(7) \cite{Carnap1950Probability}.
\end{proof}
\end{annexitem}

\begin{annexitem}{annex:p15.carnap-direct}{thm:p15-carnap-direct}{Known counts and hypergeometric direct inference}
\begin{proof}
Conditional on a fixed population count vector, all specified coarse
population words have equal probability under any symmetric measure:
a name permutation taking one to another is a measure-preserving
bijection of their full-state ranges. There are
\(N!/\prod_j n_j!\) such words. Fix a specified sample word with
counts \(u_j\). If any \(u_j>n_j\), there are no extensions.
Otherwise its remaining positions can be filled in
\((N-s)!/\prod_j(n_j-u_j)!\) ways. The conditional probability
of this specified sample word is their number divided by the number of
population words. There are \(s!/\prod_j u_j!\) specified sample
words in the sample count report. Multiplying and canceling factorials
gives \(\prod_j\binom{n_j}{u_j}/\binom Ns\).

For a one-name sample this becomes \(n_j/N\). Each of the
\(s\) named sample indicators has that mean; their sum has mean
\(sn_j/N\). To compute the fixed-\(s\) limit, write the mass as
\[
 \frac{s!}{\prod_j u_j!}
 \frac{\prod_j n_j^{\underline{u_j}}}{N^{\underline s}},
 \qquad x^{\underline d}=\prod_{i=0}^{d-1}(x-i).
\]
Divide numerator and denominator by \(N^s\). Each finite factor
\((n_j-i)/N\) tends to \(r_j\); each denominator factor
\((N-i)/N\) tends to one. This proves the limit, including zero
limiting frequencies and empty products. If a count is impossible along
some terms of a sequence, its probability is zero there and the claimed
zero limit follows for any category requiring positive \(u_j\) but
having \(r_j=0\). For the five-name binary example the count
formula gives \(3/10\). This reconstructs the direct inference of
T94-1 and the fixed-sample limiting step toward T95-1
\cite{Carnap1950Probability}.
\end{proof}
\end{annexitem}

\begin{annexitem}{annex:p15.carnap-analogy}{thm:p15-carnap-analogy}{Width-weighted two-individual analogy}
\begin{proof}
The initial probability of \(B(c)\) is \(w_B/k\), and of
\(A(c)\) is \(w_A/k\). Since \(B\subseteq A\), their
intersection is \(B(c)\). Taking the ratio proves
\(c^*(B(c),A(c))=w_B/w_A\).

The predictive theorem gives
\(c^*(B(c),B(b))=(w_B+1)/(k+1)\) and
\(c^*(A(c),B(b))=(w_A+1)/(k+1)\): the observed individual belongs
to both properties. Dividing these values, again using inclusion of
\(B(c)\) in \(A(c)\), gives \((w_B+1)/(w_A+1)\).
Subtracting \(w_B/w_A\) gives
\((w_A-w_B)/(w_A(w_A+1))\), strictly positive by the width
assumption. With widths \(1,2\), the two values are \(1/2,2/3\).
This is equation (8), sec.~110D, p.~569
\cite{Carnap1950Probability}; it makes no metric-similarity claim.
\end{proof}
\end{annexitem}

\begin{annexitem}{annex:p15.carnap-universal}{thm:p15-carnap-universal}{Universal and instance confirmations}
\begin{proof}
Use the predictive theorem on the remainder of the domain with every
future individual assigned to \(M\). The conditional probability of
each successive positive result is
\((s+w+i)/(s+k+i)\); their product for \(0\leq i<N-s\)
is the asserted rising-factorial ratio. Alternatively, the grouped
count formula has a single nonzero future count and simplifies to that
same ratio. The next-instance value is the first factor.

Put \(d=k-w>0\). For a remainder of length \(t\), the product
is \(\prod_{i=0}^{t-1}(1-d/(s+k+i))\). All factors are positive
and at most one. The inequality \(1-x\leq\exp(-x)\) for
\(x\geq0\), proved by differentiating
\(\exp(-x)-(1-x)\), bounds it above by
\(\exp(-d\sum_{i=0}^{t-1}1/(s+k+i))\).
The sum diverges: comparison with
\(\int_{s+k}^{s+k+t}dx/x\) gives a lower bound tending to infinity.
The product therefore tends to zero. For \(w=1,k=2\) it telescopes
to \((s+1)/(N+1)\). This is a finite-system limit calculation;
no existence or computation theorem for arbitrary quantified sentences
in an infinite language has been inferred from it.

The unqualified instance of \(A\to B\) is the complement of
\(A\land\neg B\), so its probability is
\(1-(s_-+w_-)/(s+k)\). To condition on the fresh individual's
being \(A\), divide the probability of \(A\land B\) by that
of the union \((A\land B)\cup(A\land\neg B)\). The common
denominator \(s+k\) cancels, giving the qualified ratio. Its
denominator is positive by the width hypotheses.
The numerical values \(101/102\), \(101/1{,}000{,}001\), and
the white-swan ratio follow by substitution. If the evidence already
contains a negative instance, its conjunction with the unrestricted law
has empty range, so regular confirmation of that law is zero. The source
distinctions are sec.~110F--H, equations (10)--(17)
\cite{Carnap1950Probability}.
\end{proof}
\end{annexitem}

\begin{annexitem}{annex:p15.carnap-history-family}{thm:p15-carnap-history-family}{A predictive family and a category-symmetric obstruction}
\begin{proof}
Start with the empty word, mass one. Given a word of length \(n\),
the positive next-category numbers \((n_i+\alpha_i)/(n+\alpha)\)
sum to one. Multiplying these numbers along a word gives
\(\prod_i\alpha_i^{\overline{n_i}}/\alpha^{\overline n}\):
the \(i\)-category factors successively use
\(\alpha_i,\ldots,\alpha_i+n_i-1\), regardless of their order.
Induction on length establishes normalization and fitting. Dependence
only on the counts establishes exchangeability. The predictive formula
is also the ratio of an extended word's mass to its prefix mass.
Splitting its numerator gives the two weighted terms in the reader.
For \(\alpha_i=1\), \(\alpha=k\), the mass is
\((k-1)!\prod_i n_i!/(n+k-1)!\), agreeing with \(m^*\).
This constructs the family; it does not prove its characterization by
the sufficientness postulate.

For the counterexample, let \(H_j\) be the model making die face
\(j\) twice as likely as each other face, and give the six models equal
prior weight. Each model assigns a specified word likelihood
\(7^{-n}2^{n_j}\). The posterior weight of \(H_j\) is therefore
\(2^{n_j}/\sum_i2^{n_i}\). Face one has next probability \(2/7\)
under \(H_1\) and \(1/7\) otherwise. Averaging gives exactly the
displayed expression. Each model is iid, so the mixture is exchangeable;
permuting category names merely permutes its six equally weighted models,
so it is category-symmetric. The two vectors in the reader have different
denominators: the first is \(2^{20}+5\cdot2^{10}\), the second
\(2^{20}+2^{30}+4\cdot2^5\). The latter is strictly larger, while
the numerator and \(n_1,n\) are the same. Thus sufficientness fails.
The historical sources for the family and example are secs.~5 and 6.4,
pp.~274--281 \cite{Zabell2011CarnapInductiveInference}.
\end{proof}
\end{annexitem}

\begin{annexitem}{annex:p15.carnap-history-finite}{thm:p15-carnap-history-finite}{Finite mixtures, nonextendability, and correlation}
\begin{proof}
Exchangeability makes every word in a count-vector orbit equally likely.
If the orbit has total probability \(v\), then each member has mass
\(v/|O|\). Hence summing the uniform orbit laws weighted by their
orbit probabilities recovers the original law. Orbits of weight zero
simply contribute zero. This is the finite representation, without an
infinite iid representation assumption.

For an exchangeable triple extending the \(RB,BR\) law, every
triple with a repeated letter would have probability zero: permute the
repeated letters to its first two positions, whose marginal \(RR\)
or \(BB\) has probability zero. Every binary triple has a repeated
letter, so all eight triple masses would be zero, contradicting
normalization. This proves nonextendability.

For \(m\geq2\) exchangeable nondegenerate Bernoulli indicators,
write their common variance as \(\sigma^2>0\) and their common
pair covariance as \(\rho\sigma^2\). Expanding the variance of
their sum gives
\[
 0\leq\operatorname{Var}\Bigl(\sum_{i=1}^mX_i\Bigr)
   =m\sigma^2+m(m-1)\rho\sigma^2.
\]
Division gives \(\rho\geq-1/(m-1)\). If arbitrarily long
exchangeable extensions preserve the pair marginal, they preserve
\(\rho\); letting their lengths grow proves \(\rho\geq0\).
For \(p=P(X_1=1)\in(0,1)\),
\[
 P(X_2=1\mid X_1=1)-p
 =\frac{P(X_1=X_2=1)-p^2}{p}
 =\frac{\operatorname{Cov}(X_1,X_2)}p\geq0.
\]
No claim that all finite exchangeable laws meet this inequality has been
made. In a known three-red/two-black urn, uniform sampling without
replacement gives \(3/5\) for the first red and \(2/4=1/2\)
for a second after the first; this illustrates negative relevance.
The source arguments are secs.~7--8
\cite{Zabell2011CarnapInductiveInference}.
\end{proof}
\end{annexitem}

\begin{annexitem}{annex:p15.carnap-history-universal}{thm:p15-carnap-history-universal}{An alternative initial allocation and uniform induction}
\begin{proof}
The all-positive law assigns mass one to the constant positive word in
each finite domain. It is normalized and fits across domains. Mixing
it with the fitting binary \(m^*\) with weights \(a,1-a\)
preserves normalization and fitting and is regular because \(1-a>0\).
For \(s\) specified positive observations, the evidence probability
is \(a+(1-a)/(s+1)\). In a domain of size \(N\geq s\), the
all-positive universal event has probability \(a+(1-a)/(N+1)\)
and implies the evidence. Its conditional probability is therefore the
ratio of these two positive quantities. Letting \(N\to\infty\)
gives the reader's positive limit. This is an independently constructed
example of the changed-prior remedy discussed in sec.~6, not a
reinterpretation of the original \(c^*\) limit.

For the second assertion put \(F_n=\bigcap_{j=1}^n E_j\) and
\(u_n=P(F_n)\). Each \(P(H\setminus E_j)=0\), so finite
subadditivity gives \(P(H\setminus F_n)=0\). Hence
\(u_n\geq P(H)>0\). Since \(F_{n+1}\subseteq F_n\),
\((u_n)\) is decreasing and has a limit
\(\ell\geq P(H)>0\). For any finite \(t\geq1\),
\[
 P(E_{n+1}\land\cdots\land E_{n+t}\mid F_n)
 =\frac{u_{n+t}}{u_n},\qquad
 0\leq1-\frac{u_{n+t}}{u_n}\leq1-\frac\ell{u_n}.
\]
The final bound tends to zero independently of \(t\), proving
uniform convergence. Only finite intersections and probability
identities were used; exchangeability was not. This is the explicit
argument in sec.~9, p.~288
\cite{Zabell2011CarnapInductiveInference}.
\end{proof}
\end{annexitem}
